% ─────────────────────────────────────────────────────────────────────────────
%  One-page, two-column résumé template
%  by Lucian Medrihan · https://1vcian.me/Portfolio
%
%  Free to use and adapt for your own résumé.
%
%  Compile with pdfLaTeX (twice, so the links settle):
%      pdflatex resume_template.tex
%  or just open it in Overleaf: it is a single file with no custom class.
%
%  How to make it yours
%    1. Search for "EDIT" and replace the placeholder text.
%    2. Change the two colours below if you want a different look.
%    3. Add or remove \experienceitem / \projectitem blocks to fit one page.
%       If the page overflows, cut content before shrinking the font.
% ─────────────────────────────────────────────────────────────────────────────

\documentclass[a4paper,10pt]{extarticle}

% Packages
\usepackage[T1]{fontenc}
\usepackage[utf8]{inputenc}
\usepackage{geometry}
\usepackage{hyperref}
\usepackage{fontawesome5}
\usepackage{titlesec}
\usepackage{enumitem}
\usepackage{xcolor}
\usepackage{tabularx}
\usepackage{multirow}
\usepackage{tikz}
\usepackage{lmodern}

% ── Colours ──────────────────────────────────────────────────────────────────
% primary   : name, section titles, timeline dots, filled stars
% linkcolor : links and the highlighted keywords
\definecolor{primary}{RGB}{11, 25, 47}
\definecolor{darkgray}{RGB}{50, 50, 50}
\definecolor{lightgray}{RGB}{245, 245, 245}
\definecolor{linkcolor}{RGB}{70, 130, 210}

% Shorthand for a highlighted keyword: \hl{Python}
\newcommand{\hl}[1]{\textcolor{linkcolor}{\textbf{#1}}}

% ── Links ────────────────────────────────────────────────────────────────────
\hypersetup{
  colorlinks=true,
  linkcolor=linkcolor,
  filecolor=linkcolor,
  urlcolor=linkcolor,
  citecolor=linkcolor,
  pdftitle={Your Name CV},     % EDIT
  pdfauthor={Your Name},       % EDIT
  pdfborder={0 0 1},
  breaklinks=true,
  bookmarksopen=true,
  bookmarksnumbered=true
}

% ── Page ─────────────────────────────────────────────────────────────────────
\geometry{
  top=0.6cm,
  bottom=0.6cm,
  left=1cm,
  right=1cm
}

\setlength{\parindent}{0pt}

% Break lines only between words: no hyphenated words, no stretched spacing
\hyphenpenalty=10000
\exhyphenpenalty=10000
\raggedright

% Section titles: bold, coloured, with a rule underneath
\titleformat{\section}
  {\large\bfseries\color{primary}}
  {}{0em}{}[\titlerule]
\titlespacing{\section}{0pt}{6pt}{3pt}

% ── Building blocks ──────────────────────────────────────────────────────────

% A small grey pill: \skill{Docker}
\newcommand{\skill}[1]{
  \colorbox{lightgray}{\footnotesize #1}
}

% The dot that starts a timeline entry
\newcommand{\timelinedot}{
  \tikz\draw[primary, fill=primary] (0,0) circle (2pt);
}

% \experienceitem{Role}{Company}{Dates}{Place}{ \item ... \item ... }
\newcommand{\experienceitem}[5]{
  \timelinedot~\textbf{#1}\\
  \textit{#2} \hfill \textit{#3}\\
  #4\\
  \begin{itemize}[leftmargin=1em,noitemsep,topsep=0pt,parsep=0pt]
    #5
  \end{itemize}
}

% \projectitem{Name}{Status}{One-line description}{website or null}{GitHub or null}
%   \item ...
%   \item ...
% \end{itemize}
%
% The macro opens the bullet list for you: write the \item lines right after it
% and close them with \end{itemize}. Pass the word null to hide an icon.
\newcommand{\projectitem}[5]{
  \textbf{#1} \hfill \textit{#2}\\
  #3 \hfill
  \ifnum\pdfstrcmp{#4}{null}=0\else\href{#4}{\faGlobe}\fi
  \ifnum\pdfstrcmp{#5}{null}=0\else~\href{#5}{\faGithub}\fi\\
  \begin{itemize}[leftmargin=1em,noitemsep,topsep=0pt,parsep=0pt]
}

% Five stars, the first #1 filled: \rating{4}
\newcommand{\rating}[1]{
  \foreach \n in {1,...,5}{
    \ifnum\n>#1
      \textcolor{lightgray}{$\star$}
    \else
      \textcolor{primary}{$\star$}
    \fi
  }
}

% One line of the Skills section: \skillrow{Label}{comma-separated items}
\newcommand{\skillrow}[2]{\textbf{#1:} #2\par\vspace{0.35em}}

\begin{document}

% ═════════════════════════════ LEFT COLUMN ═══════════════════════════════════
\begin{minipage}[t]{0.33\textwidth}
  \raggedright
  \begin{center}
    {\fontsize{21}{20}\selectfont\textbf{\textcolor{primary}{Your Name}}}\\   % EDIT
    \vspace{0.2cm}
    {\large Your Job Title}                                                    % EDIT
  \end{center}

  \section{About Me}
  % EDIT: two sentences. Lead with what you do and for whom.
  Engineer with N+ years building things that people actually use, for clients a reader will recognise. Say what you care about, in plain words.

  \section{Contact}
  % EDIT
  \begin{tabular}{@{}p{0.1\textwidth}p{0.9\textwidth}@{}}
    \faPhone & \href{tel:+00000000000}{+00 000 0000000} \\
    \faEnvelope & \href{mailto:you@example.com}{you@example.com} \\
    \faGlobe & \href{https://example.com}{example.com} \\
    \faLinkedin & \href{https://linkedin.com/in/yourname}{linkedin.com/in/yourname} \\
    \faGithub & \href{https://github.com/yourname}{github.com/yourname} \\
    \faMapMarker & City, Country \\
  \end{tabular}

  \section{Skills}
  % EDIT: one short row per area reads faster than a paragraph. Use \hl{}
  % on the few words you want a skimming reader to catch.
  \skillrow{Languages}{\hl{Language A}, Language B, Language C, Language D}
  \skillrow{Web}{\hl{Framework A}, Framework B, Framework C}
  \skillrow{AI}{\hl{Topic A}, Topic B, Topic C}
  \skillrow{Cloud}{\hl{Platform A}, Platform B, Tool C, CI/CD}
  \skillrow{Domain}{Tool A, Tool B, Tool C, Tool D}
  \skillrow{Soft skills}{leadership, mentoring, problem solving}

  \section{Education}
  % EDIT
  \timelinedot~\textbf{Degree in Your Field}\\
  \textit{Your University}

  \section{Languages}
  % EDIT: \rating takes a number from 0 to 5
  \begin{tabular}{@{}p{0.4\textwidth}p{0.6\textwidth}@{}}
    Language One & \rating{5} \\
    Language Two & \rating{4} \\
    Language Three & \rating{3} \\
  \end{tabular}

  \section{Achievements \& Activities}
  % EDIT
  \begin{itemize}[leftmargin=1.5em]
    \item \hl{Award or result}: one line on what it was and why it matters
    \item \hl{Competition or talk}, with where it took you
    \item Founder or member of \hl{a team or community}, and what it does
  \end{itemize}

\end{minipage}
\hfill
% ═════════════════════════════ RIGHT COLUMN ══════════════════════════════════
\begin{minipage}[t]{0.65\textwidth}
  \raggedright
  \section{Experiences}

  % EDIT: newest first. One line per bullet: an action and a number.
  \experienceitem{Your Current Role}{Company Name}{Mon 20XX - Present}{City, Country}{
    \item What you build and for whom, in one line
    \item A responsibility that shows scope: team size, budget, users
    \item Something you shipped on your own, with a \href{https://example.com}{link} if it is public
    \item A result with a number: cost cut from X to Y, time saved, accuracy
  }

  \vspace{0.25cm}

  \experienceitem{Your Previous Role}{Previous Company}{Mon 20XX - Mon 20XX}{City, Country}{
    \item Built N applications for a client, used by N people
    \item Refactored or migrated something: a metric improved by N\%
    \item Built internal tooling that made a recurring task N times faster
    \item Automated a manual process, saving N hours a week
    \item Introduced a new technology: what it achieved, with a number
    \item Processed a large volume of data: run time cut by N\%
  }

  \section{Projects \hfill \normalfont\small\href{https://github.com/yourname}{View more on GitHub}}

  % EDIT: the last two arguments are the website and the GitHub URL.
  % Write null for the one you do not have and its icon disappears.
  \projectitem{Project One [GitHub: $\star$ 00]}{Active}{What it is, in one line}{https://example.com}{https://github.com/yourname/project-one}
    \item What makes it technically interesting, with a number
    \item Who uses it, or how much of it there is
  \end{itemize}

  \vspace{0.2cm}

  \projectitem{Project Two}{Live}{What it is, in one line}{https://example.com}{null}
    \item The problem it solves and the result it gets
    \item The part you automated or measured
  \end{itemize}

  \vspace{0.2cm}

  \projectitem{Project Three}{Live}{What it is, in one line}{https://example.com}{null}
    \item The constraint you worked within, in numbers
    \item How you know it works: tests, users, uptime
  \end{itemize}

  \vspace{0.2cm}

  \projectitem{Project Four}{Completed}{What it is, in one line}{https://example.com}{https://github.com/yourname/project-four}
    \item What you built from scratch and its headline number
    \item Where else it ended up being used
  \end{itemize}

  \vspace{0.2cm}

  \projectitem{Project Five}{Completed}{What it is, in one line}{null}{https://github.com/yourname/project-five}
    \item One line is enough for a smaller project: what it does, with a number
  \end{itemize}

  \vspace{0.2cm}

  \projectitem{Project Six}{Completed}{What it is, in one line}{null}{https://github.com/yourname/project-six}
    \item What you made and the recognition it received
  \end{itemize}

\end{minipage}

\vfill
\begin{center}
  % EDIT: an optional closing line, or delete this block
  \small\textit{A short sign-off: where to find more, or a fun way to reach you.}
\end{center}

\end{document}
